\subsection{我们的工作}

为实现上述安全目标，我们设计并实现了 FoodShield 系统。该系统以国密算法为核心，围绕外卖订单通信场景构建了从匿名身份到可信审计的完整安全机制闭环。系统的技术结构主要包括四个层次：

\begin{itemize}[leftmargin=2em]
\item \textbf{匿名身份层}：通过基于 HMAC-SM3 的 PID 生成机制，将用户真实身份与通信身份分离，实现日常通信中的身份匿名。
\item \textbf{订单认证层}：通过基于 HMAC-SM3 的订单 Token 机制，验证通信参与方的订单合法性，防止伪造订单身份进入通信通道。
\item \textbf{安全通信层}：通过 WebSocket 实现订单级实时通信，结合 SM4-CBC 加密存储和 SM3 消息哈希，保障消息内容的保密性和可验证性。
\item \textbf{可信审计层}：通过 Merkle Tree 日志完整性验证和条件溯源机制，实现通信日志的密码学完整性证明和异常情况下的受控溯源。
\end{itemize}

具体而言，我们的主要工作如下：

\textbf{（1）基于 HMAC-SM3 的 PID 匿名身份机制。}系统在用户注册时由平台服务器基于主密钥$K_{master}$、用户真实标识$userID$和随机数$r$通过 HMAC-SM3 生成匿名身份标识 PID，即$\mathrm{PID} = \mathrm{HMAC}(K_{master},\, userID \parallel r)$。PID 在订单通信过程中替代真实身份参与认证和通信，骑手端和管理员端默认不直接展示 PID 原文，而是展示其 SM3 摘要值。该机制遵循最小知晓原则，在保障业务功能的前提下最大限度减少身份信息的暴露。

\textbf{（2）基于 HMAC-SM3 的订单 Token 认证机制。}用户创建订单后，服务器以订单号$orderID$、匿名身份标识$PID$、时间戳$timestamp$和 128 bit 随机数$nonce$作为待认证消息，使用订单认证密钥$K_{order}$计算 HMAC-SM3 值作为订单 Token。用户进入订单通信通道时需提交订单号和 Token，服务器重新计算并验证 Token 的合法性、订单归属关系和时效性。时间戳和随机数的引入有效降低了 Token 被重放或预测的风险。

\textbf{（3）基于 SM4-CBC 的消息加密存储与 SM3 哈希审计。}用户与骑手之间的通信消息经服务器处理后，使用 SM4-CBC 模式进行加密存储，管理员端默认不展示消息明文，而是展示消息哈希、时间戳、订单号等审计信息。同时，系统对每条消息日志计算绑定订单号、消息序号、发送方角色、时间戳、密文摘要和前序哈希的 SM3 摘要作为 Merkle Tree 叶子节点，为日志完整性验证奠定基础。

\textbf{（4）基于 Merkle Tree 的日志完整性验证机制。}系统采用快照机制保存订单通信日志的 Merkle Root。当订单通信结束或管理员手动触发快照时，系统对当前订单的消息日志集合计算 Merkle Root 并写入快照表。后续验证时重新计算 Root 并与历史快照对比，若不一致则说明日志被篡改。在理想哈希假设下，攻击者伪造指定 SM3 摘要的复杂度约为$2^{256}$，寻找任意碰撞的复杂度约为$2^{128}$，为日志完整性提供了较高强度的密码学保障。

\textbf{（5）条件溯源与可信审计闭环。}当发生投诉或审计需求时，用户或骑手可基于订单号提交溯源申请，管理员经权限校验后执行条件溯源流程。溯源过程以订单号为入口，经 PID 映射恢复用户真实身份，且全程记录审计日志。溯源前必须先通过 Merkle Tree 完整性验证，确保溯源所依据的通信记录未被篡改。

系统的关键技术参数如表\ref{tab:tech-params}所示。

\begin{table}[H]
\centering
\caption{FoodShield 系统关键技术参数}
\label{tab:tech-params}
\begin{tabular}{llp{7cm}}
\toprule
技术组件 & 参数 & 说明 \\
\midrule
SM3 密码杂凑 & 256 bit 输出 & 用于消息摘要、PID 哈希、Merkle 节点计算 \\
SM4 对称加密 & 128 bit 密钥 / 128 bit 分组 & CBC 模式，用于消息加密存储 \\
HMAC-SM3 & 256 bit 输出 & 用于 PID 生成和订单 Token 计算 \\
订单 Token & 含 128 bit nonce & 防重放设计，绑定订单号与匿名身份 \\
Merkle Root & 256 bit & 日志完整性标识，通过快照机制保存 \\
PID & HMAC 输出 & 匿名身份标识，与真实身份受控映射 \\
\bottomrule
\end{tabular}
\end{table}
